\documentclass[tikz,border=10pt]{standalone}
%\usepackage[fontset=lxgw]{ctex}
\usepackage{ctex}
\usepackage{xcolor}

% 定义颜色
\definecolor{arxivred}{HTML}{B31B1B} % 康奈尔红/arXiv 经典色
\definecolor{softgray}{HTML}{F5F5F5}
\definecolor{darkgray}{HTML}{333333}

\usetikzlibrary{arrows.meta, positioning, shadows}

\begin{document}
	
	\begin{tikzpicture}[
		node distance=2cm,
		every node/.style={font=\kaishu},
		timeline node/.style={
			fill=white,
			draw=arxivred,
			thick,
			rounded corners=3pt,
			inner sep=6pt,
			text width=3.5cm,
			align=center,
			drop shadow={opacity=0.2}
		},
		year node/.style={
			fill=arxivred,
			text=white,
			font=\sffamily\bfseries,
			inner sep=3pt,
			rounded corners=2pt
		}
		]
		
		% 绘制主轴线
		\draw[arxivred!30, line width=4pt, -{Stealth[length=10pt]}] (-1,0) -- (18,0);
		\draw[arxivred, line width=1.5pt] (-1,0) -- (17.5,0);
		
		% --- 1991 ---
		\node[year node] (y2023) at (0,0) {2023};
		\node[timeline node, above=0.5cm of y2023] (n1) {
			\textbf{起步：确定研究方向} \\
			\small 基于多智能体强化学习的无人艇编队控制研究
		};
		\draw[arxivred, thick] (y2023) -- (n1.south);
		
		% --- 2025 ---
		\node[year node] (y2025) at (4,0) {2025};
		\node[timeline node, below=0.5cm of y2025] (n2) {
			\textbf{开题：确定研究方向} \\
			\small 更名为 \textbf{arXiv.org}，项目移至康奈尔大学图书馆管理
		};
		\draw[arxivred, thick] (y2025) -- (n2.north);
		
		% --- 2026 ---
		\node[year node] (y2026) at (8.5,0) {2026};
		\node[timeline node, above=0.5cm of y2026] (n3) {
			\textbf{中期检查：云端迁移} \\
			\small 实验环境搭建、小论文撰写、开展实验
		};
		\draw[arxivred, thick] (y2026) -- (n3.south);
		
		% --- 2025 (Early) ---
		\node[year node] (y2027a) at (13,0) {2027.02};
		\node[timeline node, below=0.5cm of y2027a] (n4) {
			\textbf{实验：仿真场景验证} \\
			\small 全面仿真验证，ROS2仿真验证
		};
		\draw[arxivred, thick] (y2027a) -- (n4.north);
		
		% --- 2025 (Late) ---
		\node[year node] (y2025b) at (16.5,0) {2027};
		\node[timeline node, above=0.5cm of y2025b] (n5) {
			\textbf{毕业论文：答辩} \\
			\small 完善毕业论文，准备毕业答辩
		};
		\draw[arxivred, thick] (y2025b) -- (n5.south);
		
	\end{tikzpicture}
	
\end{document} 